\documentclass[11pt,dvips]{jreport}    
%\documentclass{ujarticle} 
\usepackage{amsmath,amssymb} 
%\documentclass[11pt,dvips,lines=50]{jsarticle}
\usepackage{newtxtext,newtxmath}
%\numberwithin{equation}{section}
%\usepackage[hiresbb]{graphicx}
\usepackage{graphicx}
%\usepackage{figure}
%\usepackage[dvipdfmx]{graphicx}%エラーなしで画像が出力されなかったので{graphix}にした
\usepackage{wrapfig}
\usepackage{geometry}
\geometry{left=18mm,right=10mm,top=10mm,bottom=5mm}
\setlength{\textheight}{32cm}
%\renewcommand{\bibname}{参考文献}
\usepackage[hang,small,bf]{caption}
\usepackage[subrefformat=parens]{subcaption}
\usepackage{comment}
\usepackage{ascmac}
%\usepackage{fancybx}%入ってなかった
\usepackage{caption} %キャプションの文字サイズを変えたい
\usepackage[bold]{otf}
\usepackage{tikz-feynhand}
\usepackage[misc,geometry,weather,clock]{ifsym}
\usepackage{siunitx}
\usepackage{placeins}

\captionsetup{compatibility=false}
\def\MARU#1{{\rm\ooalign{\hfil\lower.168ex\hbox{#1}\hfil \crcr\mathhexbox20D}}}




% My commands %%%%%%%%%%%%%%%%%%%%%%%%%%%%
\newcommand{\PL}[4]{
\begin{flushright}
{\LARGE\underline{\ 2026\hspace{5pt} /\hspace{1.5cm}/\hspace{1.5cm}}}\\
\vspace{10pt}
{\LARGE\underline{\ Name:  \hspace{3cm}}}\\
\end{flushright}
\vspace{-27mm}
\begin{flushleft}
\hspace{-20pt}{\large{\PaperPortrait \ Manual Check log (E71a Central Area)}}\\
\vspace{10pt}
{\LARGE\underline{\ Kdrive \ ECC#1  \  PL#2}}\\
\end{flushleft}%
\vspace{0.7cm}
#3
#4
\FreeSpace
\FloatBarrier
\newpage
}

%%event
\newcommand{\Events}[1]{
\vspace{-1cm}
\begin{center}
    			\begin{tabular}{ |p{5mm} |p{15mm}|p{26mm}|p{24mm}|p{24mm}|p{20mm}|p{40mm}|} \hline
				 &Event &film左下角に移動&オイル挿し &lens表面合わせ &移動速度調整&Comment \\
				\hline \hline
				 #1
			\end{tabular}
 	\end{center}
\FloatBarrier
\newpage
}

%% Area
\newcommand{\AreaList}[1]{%   コマンド定義
1. フィルムの位置合わせ\\
\qquad 手順：SetPlateを編集\ra film左下角に移動\ra 三か所にオイル\ra lens表面合わせ\ra 移動速度調整
 \begin{center}
    			\begin{tabular}{ |c|c|c|} \hline
				Area& Check&Comment \\
				\hline \hline
				 #1
			\end{tabular}
 	\end{center}
}%

\newcommand{\Area}[1]{
	 #1&\hspace{1cm}&\hspace{15cm}  \rule[0mm]{0mm}{0.9mm} \\ \hline
}



\newcommand{\EventList}[1]{
2. スキャン(Captureをクリックすると始まる)\\
\qquad 手順：オイル挿し\ra lensの位置合わせ\ra　移動速度調整\ra OKをクリック\\
\vspace{-0.5cm}
\begin{center}
    			\begin{tabular}[h]{ |c|c|r|c|c|} \hline%{ |p{5mm} |p{15mm}|p{26mm}|p{24mm}|p{24mm}|p{20mm}|p{40mm}|} \hline
				 &Area &Event &check&Comment\hspace{11cm} \\
				\hline \hline
				 #1
			\end{tabular}
 	\end{center}
}

\newcommand{\Event}[3]{
				#2& #3& #1 & & \rule[0mm]{0mm}{0.9mm} \\ \hline
}

\newlength{\FreeSpaceHt}
\newlength{\FreeSpaceOverhead}
\setlength{\FreeSpaceOverhead}{7cm}% itembox自体の罫線・タイトル分の余白(実測して調整済み)
\newcommand{\FreeSpace}{%
	\vspace{0.5cm}%
	\setlength{\FreeSpaceHt}{\dimexpr\pagegoal-\pagetotal-\FreeSpaceOverhead\relax}%
	\ifdim\FreeSpaceHt<0pt \setlength{\FreeSpaceHt}{0pt}\fi
	\noindent\begin{itembox}[l]{\Large{Memo}}
		\rule{0pt}{\FreeSpaceHt}
	\end{itembox}
}
\newcommand{\ra}{$\rightarrow$}
%%%%%%%%%%%%%%%%%%%%%%%%%%%%%%%%%%%%%%%%%%%%%


\begin{document}
\thispagestyle{empty}
\PL{9}{017}{
	\AreaList{
		\Area{6}
	}	
}{
	\EventList{
		\Event{14017}{o}{6}
		\Event{222530}{ }{6}
		
	}
}
\PL{9}{018}{
	\AreaList{
		\Area{2}
		\Area{3}
		\Area{4}
		\Area{6}
	}	
}{
	\EventList{
		\Event{783}{o}{2}
		\Event{2243192}{ }{2}
		\Event{3521}{o}{3}
		\Event{933259}{ }{3}
		\Event{7216}{o}{4}
		\Event{432413}{ }{4}
		\Event{14017}{ }{6}
		\Event{441985}{ }{6}
		
	}
}
\PL{9}{019}{
	\AreaList{
		\Area{1}
		\Area{2}
		\Area{3}
		\Area{4}
		\Area{6}
	}	
}{
	\EventList{
		\Event{3668}{o}{1}
		\Event{1349498}{ }{1}
		\Event{783}{ }{2}
		\Event{918079}{ }{2}
		\Event{3521}{ }{3}
		\Event{692533}{ }{3}
		\Event{7216}{ }{4}
		\Event{212769}{ }{4}
		\Event{2806}{o}{4}
		\Event{1784183}{ }{4}
		\Event{14017}{ }{6}
		\Event{423002}{ }{6}
		
	}
}
\PL{9}{020}{
	\AreaList{
		\Area{1}
		\Area{2}
		\Area{3}
		\Area{4}
		\Area{6}
	}	
}{
	\EventList{
		\Event{3668}{ }{1}
		\Event{1120386}{ }{1}
		\Event{783}{ }{2}
		\Event{697216}{ }{2}
		\Event{3521}{ }{3}
		\Event{795348}{ }{3}
		\Event{3301}{o}{4}
		\Event{3002960}{ }{4}
		\Event{7216}{ }{4}
		\Event{352201}{ }{4}
		\Event{2806}{ }{4}
		\Event{467211}{ }{4}
		\Event{6206}{o}{6}
		\Event{2280516}{ }{6}
		\Event{14017}{ }{6}
		\Event{417912}{ }{6}
		
	}
}
\PL{9}{021}{
	\AreaList{
		\Area{1}
		\Area{2}
		\Area{3}
		\Area{4}
		\Area{6}
	}	
}{
	\EventList{
		\Event{3668}{ }{1}
		\Event{262182}{ }{1}
		\Event{783}{ }{2}
		\Event{3303}{ }{2}
		\Event{3521}{ }{3}
		\Event{1507084}{ }{3}
		\Event{3301}{ }{4}
		\Event{979914}{ }{4}
		\Event{7216}{ }{4}
		\Event{2079793}{ }{4}
		\Event{2806}{ }{4}
		\Event{2336297}{ }{4}
		\Event{717}{o}{4}
		\Event{1112725}{ }{4}
		\Event{6206}{ }{6}
		\Event{543530}{ }{6}
		
	}
}
\PL{9}{022}{
	\AreaList{
		\Area{1}
		\Area{4}
		\Area{5}
		\Area{6}
	}	
}{
	\EventList{
		\Event{3668}{ }{1}
		\Event{1801066}{ }{1}
		\Event{14348}{o}{4}
		\Event{1375942}{ }{4}
		\Event{3301}{ }{4}
		\Event{1437421}{ }{4}
		\Event{2806}{ }{4}
		\Event{239223}{ }{4}
		\Event{717}{ }{4}
		\Event{32989}{ }{4}
		\Event{12366}{o}{5}
		\Event{91160}{ }{5}
		\Event{6206}{ }{6}
		\Event{406468}{ }{6}
		
	}
}
\PL{9}{023}{
	\AreaList{
		\Area{4}
		\Area{5}
		\Area{6}
	}	
}{
	\EventList{
		\Event{14348}{ }{4}
		\Event{734807}{ }{4}
		\Event{3301}{ }{4}
		\Event{251810}{ }{4}
		\Event{717}{ }{4}
		\Event{30698}{ }{4}
		\Event{12366}{ }{5}
		\Event{329377}{ }{5}
		\Event{6206}{ }{6}
		\Event{573874}{ }{6}
		
	}
}
\PL{9}{024}{
	\AreaList{
		\Area{2}
		\Area{3}
		\Area{4}
		\Area{5}
	}	
}{
	\EventList{
		\Event{8785}{o}{2}
		\Event{3151070}{ }{2}
		\Event{4468}{o}{3}
		\Event{686602}{ }{3}
		\Event{14348}{ }{4}
		\Event{1525170}{ }{4}
		\Event{717}{ }{4}
		\Event{284727}{ }{4}
		\Event{12366}{ }{5}
		\Event{1455183}{ }{5}
		
	}
}
\PL{9}{025}{
	\AreaList{
		\Area{2}
		\Area{3}
		\Area{4}
		\Area{5}
		\Area{6}
	}	
}{
	\EventList{
		\Event{8785}{ }{2}
		\Event{526480}{ }{2}
		\Event{4468}{ }{3}
		\Event{517842}{ }{3}
		\Event{14348}{ }{4}
		\Event{992994}{ }{4}
		\Event{12366}{ }{5}
		\Event{1821535}{ }{5}
		\Event{340}{o}{6}
		\Event{466291}{ }{6}
		
	}
}
\PL{9}{026}{
	\AreaList{
		\Area{2}
		\Area{3}
		\Area{4}
		\Area{5}
		\Area{6}
	}	
}{
	\EventList{
		\Event{8785}{ }{2}
		\Event{680383}{ }{2}
		\Event{4468}{ }{3}
		\Event{255828}{ }{3}
		\Event{14369}{o}{4}
		\Event{405980}{ }{4}
		\Event{7062}{o}{4}
		\Event{874888}{ }{4}
		\Event{10722}{o}{5}
		\Event{689661}{ }{5}
		\Event{340}{ }{6}
		\Event{25999}{ }{6}
		
	}
}
\PL{9}{027}{
	\AreaList{
		\Area{2}
		\Area{3}
		\Area{4}
		\Area{5}
		\Area{6}
	}	
}{
	\EventList{
		\Event{8785}{ }{2}
		\Event{1132935}{ }{2}
		\Event{4468}{ }{3}
		\Event{702909}{ }{3}
		\Event{14369}{ }{4}
		\Event{398815}{ }{4}
		\Event{7062}{ }{4}
		\Event{420593}{ }{4}
		\Event{12178}{o}{4}
		\Event{1152433}{ }{4}
		\Event{10722}{ }{5}
		\Event{903493}{ }{5}
		\Event{340}{ }{6}
		\Event{28969}{ }{6}
		
	}
}
\PL{9}{028}{
	\AreaList{
		\Area{4}
		\Area{5}
		\Area{6}
	}	
}{
	\EventList{
		\Event{14369}{ }{4}
		\Event{849878}{ }{4}
		\Event{7062}{ }{4}
		\Event{859178}{ }{4}
		\Event{12178}{ }{4}
		\Event{301023}{ }{4}
		\Event{1018}{o}{5}
		\Event{1594422}{ }{5}
		\Event{10722}{ }{5}
		\Event{249466}{ }{5}
		\Event{340}{ }{6}
		\Event{1096948}{ }{6}
		
	}
}
\PL{9}{029}{
	\AreaList{
		\Area{3}
		\Area{4}
		\Area{5}
	}	
}{
	\EventList{
		\Event{9570}{o}{3}
		\Event{848886}{ }{3}
		\Event{14369}{ }{4}
		\Event{1748908}{ }{4}
		\Event{7062}{ }{4}
		\Event{415716}{ }{4}
		\Event{12178}{ }{4}
		\Event{269605}{ }{4}
		\Event{81}{o}{4}
		\Event{866572}{ }{4}
		\Event{1018}{ }{5}
		\Event{1373243}{ }{5}
		\Event{10722}{ }{5}
		\Event{2234398}{ }{5}
		
	}
}
\PL{9}{030}{
	\AreaList{
		\Area{1}
		\Area{3}
		\Area{4}
		\Area{5}
	}	
}{
	\EventList{
		\Event{5422}{o}{1}
		\Event{2202058}{ }{1}
		\Event{6496}{o}{1}
		\Event{947301}{ }{1}
		\Event{9570}{ }{3}
		\Event{1512626}{ }{3}
		\Event{6142}{o}{4}
		\Event{1888784}{ }{4}
		\Event{12178}{ }{4}
		\Event{2089475}{ }{4}
		\Event{81}{ }{4}
		\Event{690638}{ }{4}
		\Event{1018}{ }{5}
		\Event{1593660}{ }{5}
		
	}
}
\PL{9}{031}{
	\AreaList{
		\Area{1}
		\Area{3}
		\Area{4}
		\Area{5}
	}	
}{
	\EventList{
		\Event{5422}{ }{1}
		\Event{2980414}{ }{1}
		\Event{6496}{ }{1}
		\Event{1945137}{ }{1}
		\Event{9125}{o}{3}
		\Event{2669155}{ }{3}
		\Event{9570}{ }{3}
		\Event{671597}{ }{3}
		\Event{6142}{ }{4}
		\Event{260997}{ }{4}
		\Event{81}{ }{4}
		\Event{995011}{ }{4}
		\Event{1018}{ }{5}
		\Event{2242025}{ }{5}
		
	}
}
\PL{9}{032}{
	\AreaList{
		\Area{1}
		\Area{2}
		\Area{3}
		\Area{4}
		\Area{6}
	}	
}{
	\EventList{
		\Event{5422}{ }{1}
		\Event{2093720}{ }{1}
		\Event{6496}{ }{1}
		\Event{1950064}{ }{1}
		\Event{1111}{o}{2}
		\Event{2044251}{ }{2}
		\Event{11862}{o}{3}
		\Event{1437146}{ }{3}
		\Event{9125}{ }{3}
		\Event{1140712}{ }{3}
		\Event{9570}{ }{3}
		\Event{1376142}{ }{3}
		\Event{6142}{ }{4}
		\Event{259951}{ }{4}
		\Event{81}{ }{4}
		\Event{769537}{ }{4}
		\Event{4478}{o}{6}
		\Event{1502787}{ }{6}
		
	}
}
\PL{9}{033}{
	\AreaList{
		\Area{1}
		\Area{2}
		\Area{3}
		\Area{4}
		\Area{6}
	}	
}{
	\EventList{
		\Event{5422}{ }{1}
		\Event{1927733}{ }{1}
		\Event{6496}{ }{1}
		\Event{919340}{ }{1}
		\Event{1111}{ }{2}
		\Event{1472357}{ }{2}
		\Event{11862}{ }{3}
		\Event{157514}{ }{3}
		\Event{9125}{ }{3}
		\Event{279451}{ }{3}
		\Event{6142}{ }{4}
		\Event{705451}{ }{4}
		\Event{897}{o}{4}
		\Event{491838}{ }{4}
		\Event{4478}{ }{6}
		\Event{784250}{ }{6}
		
	}
}
\PL{9}{034}{
	\AreaList{
		\Area{2}
		\Area{3}
		\Area{4}
		\Area{6}
	}	
}{
	\EventList{
		\Event{1111}{ }{2}
		\Event{908855}{ }{2}
		\Event{11862}{ }{3}
		\Event{949276}{ }{3}
		\Event{1747}{o}{3}
		\Event{1559619}{ }{3}
		\Event{7141}{o}{3}
		\Event{1788044}{ }{3}
		\Event{9125}{ }{3}
		\Event{1789754}{ }{3}
		\Event{897}{ }{4}
		\Event{11018}{ }{4}
		\Event{4478}{ }{6}
		\Event{842543}{ }{6}
		
	}
}
\PL{9}{035}{
	\AreaList{
		\Area{2}
		\Area{3}
		\Area{4}
		\Area{6}
	}	
}{
	\EventList{
		\Event{1111}{ }{2}
		\Event{742907}{ }{2}
		\Event{11862}{ }{3}
		\Event{1057296}{ }{3}
		\Event{7141}{ }{3}
		\Event{1575451}{ }{3}
		\Event{1747}{ }{4}
		\Event{259830}{ }{4}
		\Event{897}{ }{4}
		\Event{417037}{ }{4}
		\Event{4478}{ }{6}
		\Event{214772}{ }{6}
		\Event{9056}{o}{6}
		\Event{477065}{ }{6}
		
	}
}
\PL{9}{036}{
	\AreaList{
		\Area{3}
		\Area{4}
		\Area{6}
	}	
}{
	\EventList{
		\Event{1747}{ }{3}
		\Event{503646}{ }{3}
		\Event{7141}{ }{3}
		\Event{1713177}{ }{3}
		\Event{897}{ }{4}
		\Event{11900}{ }{4}
		\Event{2306}{o}{6}
		\Event{1394318}{ }{6}
		\Event{9056}{ }{6}
		\Event{1354546}{ }{6}
		
	}
}
\PL{9}{037}{
	\AreaList{
		\Area{3}
		\Area{4}
		\Area{6}
	}	
}{
	\EventList{
		\Event{7141}{ }{3}
		\Event{845301}{ }{3}
		\Event{1747}{ }{4}
		\Event{271644}{ }{4}
		\Event{2306}{ }{6}
		\Event{934740}{ }{6}
		\Event{9056}{ }{6}
		\Event{1623876}{ }{6}
		
	}
}
\PL{9}{038}{
	\AreaList{
		\Area{2}
		\Area{4}
		\Area{5}
		\Area{6}
	}	
}{
	\EventList{
		\Event{13413}{o}{2}
		\Event{4837017}{ }{2}
		\Event{9023}{o}{4}
		\Event{1920319}{ }{4}
		\Event{13268}{o}{5}
		\Event{1261157}{ }{5}
		\Event{2306}{ }{6}
		\Event{279917}{ }{6}
		\Event{9056}{ }{6}
		\Event{276742}{ }{6}
		
	}
}
\PL{9}{039}{
	\AreaList{
		\Area{2}
		\Area{4}
		\Area{5}
		\Area{6}
	}	
}{
	\EventList{
		\Event{13413}{ }{2}
		\Event{1495609}{ }{2}
		\Event{9023}{ }{4}
		\Event{350811}{ }{4}
		\Event{12067}{o}{4}
		\Event{1718785}{ }{4}
		\Event{13268}{ }{5}
		\Event{1034345}{ }{5}
		\Event{2306}{ }{6}
		\Event{269867}{ }{6}
		
	}
}
\PL{9}{040}{
	\AreaList{
		\Area{2}
		\Area{4}
		\Area{5}
	}	
}{
	\EventList{
		\Event{13413}{ }{2}
		\Event{1389666}{ }{2}
		\Event{9023}{ }{4}
		\Event{1257422}{ }{4}
		\Event{12067}{ }{4}
		\Event{1765300}{ }{4}
		\Event{13268}{ }{5}
		\Event{146540}{ }{5}
		
	}
}
\PL{9}{041}{
	\AreaList{
		\Area{2}
		\Area{3}
		\Area{4}
		\Area{5}
	}	
}{
	\EventList{
		\Event{13413}{ }{2}
		\Event{749098}{ }{2}
		\Event{12860}{o}{3}
		\Event{3284345}{ }{3}
		\Event{9023}{ }{4}
		\Event{1427207}{ }{4}
		\Event{12067}{ }{4}
		\Event{661825}{ }{4}
		\Event{13268}{ }{5}
		\Event{1179861}{ }{5}
		
	}
}
\PL{9}{042}{
	\AreaList{
		\Area{3}
		\Area{4}
		\Area{6}
	}	
}{
	\EventList{
		\Event{15140}{o}{3}
		\Event{607077}{ }{3}
		\Event{12860}{ }{3}
		\Event{595641}{ }{3}
		\Event{12067}{ }{4}
		\Event{857902}{ }{4}
		\Event{11677}{o}{6}
		\Event{1867059}{ }{6}
		
	}
}
\PL{9}{043}{
	\AreaList{
		\Area{3}
		\Area{6}
	}	
}{
	\EventList{
		\Event{15140}{ }{3}
		\Event{1862240}{ }{3}
		\Event{12860}{ }{3}
		\Event{552486}{ }{3}
		\Event{11677}{ }{6}
		\Event{166940}{ }{6}
		
	}
}
\PL{9}{044}{
	\AreaList{
		\Area{3}
		\Area{4}
		\Area{6}
	}	
}{
	\EventList{
		\Event{15140}{ }{3}
		\Event{1874580}{ }{3}
		\Event{12860}{ }{3}
		\Event{546189}{ }{3}
		\Event{14299}{o}{4}
		\Event{567176}{ }{4}
		\Event{8067}{o}{4}
		\Event{537245}{ }{4}
		\Event{11677}{ }{6}
		\Event{150219}{ }{6}
		
	}
}
\PL{9}{045}{
	\AreaList{
		\Area{3}
		\Area{4}
		\Area{6}
	}	
}{
	\EventList{
		\Event{15140}{ }{3}
		\Event{1014256}{ }{3}
		\Event{14299}{ }{4}
		\Event{864015}{ }{4}
		\Event{8067}{ }{4}
		\Event{1446616}{ }{4}
		\Event{11677}{ }{6}
		\Event{1900107}{ }{6}
		
	}
}
\PL{9}{046}{
	\AreaList{
		\Area{2}
		\Area{3}
		\Area{4}
		\Area{6}
	}	
}{
	\EventList{
		\Event{13699}{o}{2}
		\Event{249910}{ }{2}
		\Event{8003}{o}{3}
		\Event{1393500}{ }{3}
		\Event{14299}{ }{4}
		\Event{423824}{ }{4}
		\Event{1556}{o}{4}
		\Event{943228}{ }{4}
		\Event{5517}{o}{4}
		\Event{1764945}{ }{4}
		\Event{8067}{ }{4}
		\Event{238453}{ }{4}
		\Event{7183}{o}{6}
		\Event{1143429}{ }{6}
		
	}
}
\PL{9}{047}{
	\AreaList{
		\Area{2}
		\Area{3}
		\Area{4}
		\Area{6}
	}	
}{
	\EventList{
		\Event{13699}{ }{2}
		\Event{243697}{ }{2}
		\Event{8003}{ }{3}
		\Event{3157846}{ }{3}
		\Event{14299}{ }{4}
		\Event{616929}{ }{4}
		\Event{1556}{ }{4}
		\Event{263561}{ }{4}
		\Event{5517}{ }{4}
		\Event{1749952}{ }{4}
		\Event{8067}{ }{4}
		\Event{975990}{ }{4}
		\Event{7183}{ }{6}
		\Event{723883}{ }{6}
		
	}
}
\PL{9}{048}{
	\AreaList{
		\Area{2}
		\Area{3}
		\Area{4}
		\Area{6}
	}	
}{
	\EventList{
		\Event{13699}{ }{2}
		\Event{233722}{ }{2}
		\Event{8003}{ }{3}
		\Event{1922005}{ }{3}
		\Event{1556}{ }{4}
		\Event{1215666}{ }{4}
		\Event{5517}{ }{4}
		\Event{178750}{ }{4}
		\Event{7183}{ }{6}
		\Event{677501}{ }{6}
		
	}
}
\PL{9}{049}{
	\AreaList{
		\Area{2}
		\Area{3}
		\Area{4}
		\Area{5}
		\Area{6}
	}	
}{
	\EventList{
		\Event{13699}{ }{2}
		\Event{269979}{ }{2}
		\Event{12306}{o}{3}
		\Event{1315199}{ }{3}
		\Event{1556}{ }{4}
		\Event{1248656}{ }{4}
		\Event{5517}{ }{4}
		\Event{194141}{ }{4}
		\Event{8003}{ }{5}
		\Event{638030}{ }{5}
		\Event{7183}{ }{6}
		\Event{262267}{ }{6}
		
	}
}
\PL{9}{050}{
	\AreaList{
		\Area{3}
		\Area{5}
	}	
}{
	\EventList{
		\Event{13221}{o}{3}
		\Event{1201670}{ }{3}
		\Event{9821}{o}{3}
		\Event{480035}{ }{3}
		\Event{9891}{o}{3}
		\Event{397310}{ }{3}
		\Event{12306}{ }{3}
		\Event{434347}{ }{3}
		\Event{5987}{o}{5}
		\Event{347125}{ }{5}
		
	}
}
\PL{9}{051}{
	\AreaList{
		\Area{3}
		\Area{5}
	}	
}{
	\EventList{
		\Event{13221}{ }{3}
		\Event{295699}{ }{3}
		\Event{9821}{ }{3}
		\Event{269766}{ }{3}
		\Event{9891}{ }{3}
		\Event{350631}{ }{3}
		\Event{12306}{ }{3}
		\Event{1097655}{ }{3}
		\Event{5987}{ }{5}
		\Event{778191}{ }{5}
		
	}
}
\PL{9}{052}{
	\AreaList{
		\Area{2}
		\Area{3}
		\Area{4}
		\Area{5}
		\Area{6}
	}	
}{
	\EventList{
		\Event{9491}{o}{2}
		\Event{1167424}{ }{2}
		\Event{13221}{ }{3}
		\Event{951843}{ }{3}
		\Event{9821}{ }{3}
		\Event{927795}{ }{3}
		\Event{9891}{ }{3}
		\Event{807773}{ }{3}
		\Event{12306}{ }{3}
		\Event{887034}{ }{3}
		\Event{12535}{o}{4}
		\Event{3475689}{ }{4}
		\Event{5987}{ }{5}
		\Event{176383}{ }{5}
		\Event{5678}{o}{6}
		\Event{2869400}{ }{6}
		
	}
}
\PL{9}{053}{
	\AreaList{
		\Area{2}
		\Area{3}
		\Area{4}
		\Area{5}
		\Area{6}
	}	
}{
	\EventList{
		\Event{9491}{ }{2}
		\Event{1144778}{ }{2}
		\Event{13221}{ }{3}
		\Event{1853850}{ }{3}
		\Event{9821}{ }{3}
		\Event{710273}{ }{3}
		\Event{9891}{ }{3}
		\Event{1410605}{ }{3}
		\Event{12535}{ }{4}
		\Event{250676}{ }{4}
		\Event{5987}{ }{5}
		\Event{202447}{ }{5}
		\Event{5678}{ }{6}
		\Event{440452}{ }{6}
		
	}
}
\PL{9}{054}{
	\AreaList{
		\Area{2}
		\Area{3}
		\Area{4}
		\Area{6}
	}	
}{
	\EventList{
		\Event{9491}{ }{2}
		\Event{559154}{ }{2}
		\Event{11300}{o}{3}
		\Event{1200523}{ }{3}
		\Event{4126}{o}{3}
		\Event{606244}{ }{3}
		\Event{12535}{ }{4}
		\Event{289733}{ }{4}
		\Event{5678}{ }{6}
		\Event{440715}{ }{6}
		
	}
}
\PL{9}{055}{
	\AreaList{
		\Area{2}
		\Area{3}
		\Area{4}
		\Area{6}
	}	
}{
	\EventList{
		\Event{9491}{ }{2}
		\Event{738442}{ }{2}
		\Event{11300}{ }{3}
		\Event{677686}{ }{3}
		\Event{4126}{ }{3}
		\Event{328362}{ }{3}
		\Event{12535}{ }{4}
		\Event{470777}{ }{4}
		\Event{5678}{ }{6}
		\Event{1103082}{ }{6}
		
	}
}
\PL{9}{056}{
	\AreaList{
		\Area{3}
		\Area{6}
	}	
}{
	\EventList{
		\Event{11300}{ }{3}
		\Event{549364}{ }{3}
		\Event{4126}{ }{3}
		\Event{372808}{ }{3}
		\Event{5828}{o}{6}
		\Event{722541}{ }{6}
		
	}
}
\PL{9}{057}{
	\AreaList{
		\Area{1}
		\Area{3}
		\Area{6}
	}	
}{
	\EventList{
		\Event{3069}{o}{1}
		\Event{845268}{ }{1}
		\Event{11300}{ }{3}
		\Event{473898}{ }{3}
		\Event{4126}{ }{3}
		\Event{1455621}{ }{3}
		\Event{5828}{ }{6}
		\Event{694147}{ }{6}
		
	}
}
\PL{9}{058}{
	\AreaList{
		\Area{1}
		\Area{4}
		\Area{6}
	}	
}{
	\EventList{
		\Event{3069}{ }{1}
		\Event{1483115}{ }{1}
		\Event{2209}{o}{4}
		\Event{629383}{ }{4}
		\Event{5828}{ }{6}
		\Event{39813}{ }{6}
		
	}
}
\PL{9}{059}{
	\AreaList{
		\Area{1}
		\Area{4}
		\Area{6}
	}	
}{
	\EventList{
		\Event{3069}{ }{1}
		\Event{1396769}{ }{1}
		\Event{2209}{ }{4}
		\Event{648056}{ }{4}
		\Event{5828}{ }{6}
		\Event{53771}{ }{6}
		\Event{15214}{o}{6}
		\Event{426837}{ }{6}
		
	}
}
\PL{9}{060}{
	\AreaList{
		\Area{1}
		\Area{4}
		\Area{5}
		\Area{6}
	}	
}{
	\EventList{
		\Event{3069}{ }{1}
		\Event{1715964}{ }{1}
		\Event{2209}{ }{4}
		\Event{864116}{ }{4}
		\Event{8485}{o}{5}
		\Event{1423543}{ }{5}
		\Event{15214}{ }{6}
		\Event{216486}{ }{6}
		
	}
}
\PL{9}{061}{
	\AreaList{
		\Area{4}
		\Area{5}
		\Area{6}
	}	
}{
	\EventList{
		\Event{2209}{ }{4}
		\Event{864390}{ }{4}
		\Event{8485}{ }{5}
		\Event{327172}{ }{5}
		\Event{15214}{ }{6}
		\Event{808704}{ }{6}
		
	}
}
\PL{9}{062}{
	\AreaList{
		\Area{1}
		\Area{3}
		\Area{4}
		\Area{5}
		\Area{6}
	}	
}{
	\EventList{
		\Event{10776}{o}{1}
		\Event{1451206}{ }{1}
		\Event{2243}{o}{3}
		\Event{1278666}{ }{3}
		\Event{13480}{o}{4}
		\Event{1282578}{ }{4}
		\Event{8485}{ }{5}
		\Event{323636}{ }{5}
		\Event{1592}{o}{6}
		\Event{333337}{ }{6}
		\Event{15214}{ }{6}
		\Event{811996}{ }{6}
		
	}
}
\PL{9}{063}{
	\AreaList{
		\Area{1}
		\Area{3}
		\Area{4}
		\Area{5}
		\Area{6}
	}	
}{
	\EventList{
		\Event{10776}{ }{1}
		\Event{496784}{ }{1}
		\Event{2243}{ }{3}
		\Event{1277603}{ }{3}
		\Event{484}{o}{3}
		\Event{525306}{ }{3}
		\Event{13480}{ }{4}
		\Event{1057472}{ }{4}
		\Event{8485}{ }{5}
		\Event{190890}{ }{5}
		\Event{576}{o}{5}
		\Event{1709330}{ }{5}
		\Event{1592}{ }{6}
		\Event{582454}{ }{6}
		
	}
}
\PL{9}{064}{
	\AreaList{
		\Area{1}
		\Area{2}
		\Area{3}
		\Area{4}
		\Area{5}
		\Area{6}
	}	
}{
	\EventList{
		\Event{10776}{ }{1}
		\Event{981149}{ }{1}
		\Event{5021}{o}{2}
		\Event{1997168}{ }{2}
		\Event{2243}{ }{3}
		\Event{1694583}{ }{3}
		\Event{9452}{o}{3}
		\Event{1494123}{ }{3}
		\Event{484}{ }{3}
		\Event{293941}{ }{3}
		\Event{13480}{ }{4}
		\Event{383627}{ }{4}
		\Event{12779}{o}{5}
		\Event{717662}{ }{5}
		\Event{576}{ }{5}
		\Event{152254}{ }{5}
		\Event{1592}{ }{6}
		\Event{817457}{ }{6}
		
	}
}
\PL{9}{065}{
	\AreaList{
		\Area{1}
		\Area{3}
		\Area{4}
		\Area{5}
		\Area{6}
	}	
}{
	\EventList{
		\Event{10776}{ }{1}
		\Event{532214}{ }{1}
		\Event{12779}{ }{3}
		\Event{1938545}{ }{3}
		\Event{2243}{ }{3}
		\Event{366550}{ }{3}
		\Event{9452}{ }{3}
		\Event{709294}{ }{3}
		\Event{484}{ }{3}
		\Event{1774115}{ }{3}
		\Event{13480}{ }{4}
		\Event{826243}{ }{4}
		\Event{5021}{ }{4}
		\Event{621950}{ }{4}
		\Event{2320}{o}{4}
		\Event{235117}{ }{4}
		\Event{576}{ }{5}
		\Event{116710}{ }{5}
		\Event{1592}{ }{6}
		\Event{777826}{ }{6}
		
	}
}
\PL{9}{066}{
	\AreaList{
		\Area{2}
		\Area{3}
		\Area{4}
	}	
}{
	\EventList{
		\Event{5021}{ }{2}
		\Event{811485}{ }{2}
		\Event{12779}{ }{3}
		\Event{1501295}{ }{3}
		\Event{9452}{ }{3}
		\Event{100537}{ }{3}
		\Event{484}{ }{3}
		\Event{2085315}{ }{3}
		\Event{576}{ }{3}
		\Event{788596}{ }{3}
		\Event{8190}{o}{4}
		\Event{928179}{ }{4}
		\Event{2320}{ }{4}
		\Event{11778}{ }{4}
		
	}
}
\PL{9}{067}{
	\AreaList{
		\Area{2}
		\Area{3}
		\Area{4}
		\Area{6}
	}	
}{
	\EventList{
		\Event{5021}{ }{2}
		\Event{2072630}{ }{2}
		\Event{12779}{ }{3}
		\Event{1237344}{ }{3}
		\Event{9452}{ }{3}
		\Event{749246}{ }{3}
		\Event{8190}{ }{4}
		\Event{673304}{ }{4}
		\Event{2320}{ }{4}
		\Event{11850}{ }{4}
		\Event{4089}{o}{6}
		\Event{885012}{ }{6}
		
	}
}
\PL{9}{068}{
	\AreaList{
		\Area{4}
		\Area{6}
	}	
}{
	\EventList{
		\Event{8190}{ }{4}
		\Event{327817}{ }{4}
		\Event{2320}{ }{4}
		\Event{425266}{ }{4}
		\Event{4089}{ }{6}
		\Event{651490}{ }{6}
		
	}
}
\PL{9}{069}{
	\AreaList{
		\Area{4}
		\Area{6}
	}	
}{
	\EventList{
		\Event{8190}{ }{4}
		\Event{788395}{ }{4}
		\Event{4089}{ }{6}
		\Event{1360}{ }{6}
		
	}
}
\PL{9}{070}{
	\AreaList{
		\Area{2}
		\Area{3}
		\Area{4}
		\Area{6}
	}	
}{
	\EventList{
		\Event{9858}{o}{2}
		\Event{6457849}{ }{2}
		\Event{3121}{o}{3}
		\Event{1262614}{ }{3}
		\Event{5827}{o}{4}
		\Event{2612653}{ }{4}
		\Event{4089}{ }{6}
		\Event{646443}{ }{6}
		
	}
}
\PL{9}{071}{
	\AreaList{
		\Area{2}
		\Area{3}
		\Area{4}
	}	
}{
	\EventList{
		\Event{9858}{ }{2}
		\Event{668967}{ }{2}
		\Event{3121}{ }{3}
		\Event{336908}{ }{3}
		\Event{5827}{ }{4}
		\Event{856580}{ }{4}
		
	}
}
\PL{9}{072}{
	\AreaList{
		\Area{1}
		\Area{2}
		\Area{3}
		\Area{4}
	}	
}{
	\EventList{
		\Event{1623}{o}{1}
		\Event{307429}{ }{1}
		\Event{5800}{o}{1}
		\Event{269945}{ }{1}
		\Event{1338}{o}{2}
		\Event{80282}{ }{2}
		\Event{9858}{ }{2}
		\Event{514193}{ }{2}
		\Event{3121}{ }{3}
		\Event{1087033}{ }{3}
		\Event{5827}{ }{4}
		\Event{1096149}{ }{4}
		
	}
}
\PL{9}{073}{
	\AreaList{
		\Area{1}
		\Area{2}
		\Area{3}
		\Area{4}
	}	
}{
	\EventList{
		\Event{1623}{ }{1}
		\Event{1007272}{ }{1}
		\Event{5800}{ }{1}
		\Event{774397}{ }{1}
		\Event{1338}{ }{2}
		\Event{1487316}{ }{2}
		\Event{9858}{ }{2}
		\Event{807636}{ }{2}
		\Event{3121}{ }{3}
		\Event{817243}{ }{3}
		\Event{5827}{ }{4}
		\Event{644894}{ }{4}
		
	}
}
\PL{9}{074}{
	\AreaList{
		\Area{1}
		\Area{2}
		\Area{3}
		\Area{5}
	}	
}{
	\EventList{
		\Event{1623}{ }{1}
		\Event{859193}{ }{1}
		\Event{5800}{ }{1}
		\Event{710119}{ }{1}
		\Event{1338}{ }{2}
		\Event{679164}{ }{2}
		\Event{3253}{o}{3}
		\Event{221007}{ }{3}
		\Event{14021}{o}{5}
		\Event{609539}{ }{5}
		
	}
}
\PL{9}{075}{
	\AreaList{
		\Area{1}
		\Area{2}
		\Area{3}
		\Area{5}
	}	
}{
	\EventList{
		\Event{1623}{ }{1}
		\Event{1070743}{ }{1}
		\Event{5800}{ }{1}
		\Event{905677}{ }{1}
		\Event{1338}{ }{2}
		\Event{1201229}{ }{2}
		\Event{11833}{o}{2}
		\Event{2699018}{ }{2}
		\Event{14734}{o}{2}
		\Event{719648}{ }{2}
		\Event{3253}{ }{3}
		\Event{854}{ }{3}
		\Event{14021}{ }{5}
		\Event{593255}{ }{5}
		
	}
}
\PL{9}{076}{
	\AreaList{
		\Area{2}
		\Area{3}
		\Area{5}
		\Area{6}
	}	
}{
	\EventList{
		\Event{11833}{ }{2}
		\Event{1870911}{ }{2}
		\Event{14734}{ }{2}
		\Event{1569540}{ }{2}
		\Event{3253}{ }{3}
		\Event{662408}{ }{3}
		\Event{14021}{ }{5}
		\Event{1748737}{ }{5}
		\Event{10233}{o}{6}
		\Event{1936988}{ }{6}
		
	}
}
\PL{9}{077}{
	\AreaList{
		\Area{2}
		\Area{3}
		\Area{5}
		\Area{6}
	}	
}{
	\EventList{
		\Event{11833}{ }{2}
		\Event{301970}{ }{2}
		\Event{14734}{ }{2}
		\Event{13005}{ }{2}
		\Event{3253}{ }{3}
		\Event{882964}{ }{3}
		\Event{14021}{ }{5}
		\Event{406243}{ }{5}
		\Event{10233}{ }{6}
		\Event{355522}{ }{6}
		
	}
}
\PL{9}{078}{
	\AreaList{
		\Area{2}
		\Area{6}
	}	
}{
	\EventList{
		\Event{11833}{ }{2}
		\Event{783645}{ }{2}
		\Event{14734}{ }{2}
		\Event{248360}{ }{2}
		\Event{10233}{ }{6}
		\Event{376882}{ }{6}
		
	}
}
\PL{9}{079}{
	\AreaList{
		\Area{6}
	}	
}{
	\EventList{
		\Event{10233}{ }{6}
		\Event{141046}{ }{6}
		
	}
}
\PL{9}{080}{
	\AreaList{
		\Area{1}
	}	
}{
	\EventList{
		\Event{6875}{o}{1}
		\Event{1768656}{ }{1}
		
	}
}
\PL{9}{081}{
	\AreaList{
		\Area{1}
	}	
}{
	\EventList{
		\Event{6875}{ }{1}
		\Event{464800}{ }{1}
		
	}
}
\PL{9}{082}{
	\AreaList{
		\Area{1}
	}	
}{
	\EventList{
		\Event{6875}{ }{1}
		\Event{498293}{ }{1}
		
	}
}
\PL{9}{083}{
	\AreaList{
		\Area{1}
	}	
}{
	\EventList{
		\Event{6875}{ }{1}
		\Event{980052}{ }{1}
		
	}
}
\PL{9}{084}{
	\AreaList{
		\Area{1}
		\Area{3}
		\Area{4}
	}	
}{
	\EventList{
		\Event{3743}{o}{1}
		\Event{1336241}{ }{1}
		\Event{14553}{o}{3}
		\Event{658670}{ }{3}
		\Event{4948}{o}{4}
		\Event{1352158}{ }{4}
		\Event{6374}{o}{4}
		\Event{1074090}{ }{4}
		
	}
}
\PL{9}{085}{
	\AreaList{
		\Area{1}
		\Area{3}
		\Area{4}
		\Area{5}
	}	
}{
	\EventList{
		\Event{3743}{ }{1}
		\Event{973649}{ }{1}
		\Event{14553}{ }{3}
		\Event{316230}{ }{3}
		\Event{4948}{ }{4}
		\Event{773305}{ }{4}
		\Event{6374}{ }{4}
		\Event{411008}{ }{4}
		\Event{10957}{o}{4}
		\Event{1660368}{ }{4}
		\Event{2720}{o}{5}
		\Event{267426}{ }{5}
		
	}
}
\PL{9}{086}{
	\AreaList{
		\Area{1}
		\Area{3}
		\Area{4}
		\Area{5}
		\Area{6}
	}	
}{
	\EventList{
		\Event{3743}{ }{1}
		\Event{1081122}{ }{1}
		\Event{14553}{ }{3}
		\Event{341890}{ }{3}
		\Event{6374}{ }{4}
		\Event{5122}{ }{4}
		\Event{10957}{ }{4}
		\Event{753985}{ }{4}
		\Event{6755}{o}{5}
		\Event{2636235}{ }{5}
		\Event{2720}{ }{5}
		\Event{803754}{ }{5}
		\Event{1435}{o}{6}
		\Event{851758}{ }{6}
		\Event{4948}{ }{6}
		\Event{117864}{ }{6}
		
	}
}
\PL{9}{087}{
	\AreaList{
		\Area{1}
		\Area{3}
		\Area{4}
		\Area{5}
		\Area{6}
	}	
}{
	\EventList{
		\Event{3743}{ }{1}
		\Event{860299}{ }{1}
		\Event{14553}{ }{3}
		\Event{1187319}{ }{3}
		\Event{4948}{ }{4}
		\Event{1211341}{ }{4}
		\Event{6374}{ }{4}
		\Event{415187}{ }{4}
		\Event{10957}{ }{4}
		\Event{927799}{ }{4}
		\Event{6755}{ }{5}
		\Event{2416948}{ }{5}
		\Event{2720}{ }{5}
		\Event{1231367}{ }{5}
		\Event{1435}{ }{6}
		\Event{1773020}{ }{6}
		
	}
}
\PL{9}{088}{
	\AreaList{
		\Area{4}
		\Area{5}
		\Area{6}
	}	
}{
	\EventList{
		\Event{6637}{o}{4}
		\Event{316627}{ }{4}
		\Event{10957}{ }{4}
		\Event{301354}{ }{4}
		\Event{6755}{ }{5}
		\Event{409629}{ }{5}
		\Event{2720}{ }{5}
		\Event{98215}{ }{5}
		\Event{1435}{ }{6}
		\Event{409054}{ }{6}
		
	}
}
\PL{9}{089}{
	\AreaList{
		\Area{4}
		\Area{5}
		\Area{6}
	}	
}{
	\EventList{
		\Event{6637}{ }{4}
		\Event{193221}{ }{4}
		\Event{6755}{ }{5}
		\Event{1284493}{ }{5}
		\Event{1435}{ }{6}
		\Event{231766}{ }{6}
		\Event{10182}{o}{6}
		\Event{2014509}{ }{6}
		
	}
}
\PL{9}{090}{
	\AreaList{
		\Area{2}
		\Area{4}
		\Area{6}
	}	
}{
	\EventList{
		\Event{11777}{o}{2}
		\Event{1105935}{ }{2}
		\Event{6637}{ }{4}
		\Event{988285}{ }{4}
		\Event{10182}{ }{6}
		\Event{169098}{ }{6}
		
	}
}
\PL{9}{091}{
	\AreaList{
		\Area{2}
		\Area{4}
		\Area{6}
	}	
}{
	\EventList{
		\Event{11777}{ }{2}
		\Event{2460866}{ }{2}
		\Event{6637}{ }{4}
		\Event{197717}{ }{4}
		\Event{10182}{ }{6}
		\Event{1741847}{ }{6}
		
	}
}
\PL{9}{092}{
	\AreaList{
		\Area{2}
		\Area{3}
		\Area{5}
		\Area{6}
	}	
}{
	\EventList{
		\Event{11777}{ }{2}
		\Event{1328044}{ }{2}
		\Event{10356}{o}{3}
		\Event{1438231}{ }{3}
		\Event{12692}{o}{3}
		\Event{514331}{ }{3}
		\Event{6971}{o}{5}
		\Event{514330}{ }{5}
		\Event{10182}{ }{6}
		\Event{863983}{ }{6}
		
	}
}
\PL{9}{093}{
	\AreaList{
		\Area{1}
		\Area{2}
		\Area{3}
		\Area{5}
	}	
}{
	\EventList{
		\Event{1913}{o}{1}
		\Event{809859}{ }{1}
		\Event{11777}{ }{2}
		\Event{483345}{ }{2}
		\Event{10356}{ }{3}
		\Event{1781026}{ }{3}
		\Event{12692}{ }{3}
		\Event{673202}{ }{3}
		\Event{6971}{ }{5}
		\Event{276200}{ }{5}
		
	}
}
\PL{9}{094}{
	\AreaList{
		\Area{1}
		\Area{3}
		\Area{5}
		\Area{6}
	}	
}{
	\EventList{
		\Event{4796}{o}{1}
		\Event{753200}{ }{1}
		\Event{1913}{ }{1}
		\Event{1847153}{ }{1}
		\Event{10356}{ }{3}
		\Event{514913}{ }{3}
		\Event{12692}{ }{3}
		\Event{304663}{ }{3}
		\Event{7621}{o}{3}
		\Event{603306}{ }{3}
		\Event{6971}{ }{5}
		\Event{1198846}{ }{5}
		\Event{14261}{o}{6}
		\Event{1276039}{ }{6}
		
	}
}
\PL{9}{095}{
	\AreaList{
		\Area{1}
		\Area{3}
		\Area{4}
		\Area{5}
		\Area{6}
	}	
}{
	\EventList{
		\Event{4796}{ }{1}
		\Event{510326}{ }{1}
		\Event{1913}{ }{1}
		\Event{1537141}{ }{1}
		\Event{10356}{ }{3}
		\Event{487311}{ }{3}
		\Event{12692}{ }{3}
		\Event{282562}{ }{3}
		\Event{7621}{ }{3}
		\Event{368090}{ }{3}
		\Event{14349}{o}{4}
		\Event{763431}{ }{4}
		\Event{6971}{ }{5}
		\Event{476115}{ }{5}
		\Event{14261}{ }{6}
		\Event{188929}{ }{6}
		
	}
}
\PL{9}{096}{
	\AreaList{
		\Area{1}
		\Area{3}
		\Area{4}
		\Area{5}
		\Area{6}
	}	
}{
	\EventList{
		\Event{4796}{ }{1}
		\Event{309894}{ }{1}
		\Event{1913}{ }{1}
		\Event{505122}{ }{1}
		\Event{7621}{ }{3}
		\Event{389081}{ }{3}
		\Event{14349}{ }{4}
		\Event{1176527}{ }{4}
		\Event{12675}{o}{5}
		\Event{1068694}{ }{5}
		\Event{14261}{ }{6}
		\Event{525239}{ }{6}
		
	}
}
\PL{9}{097}{
	\AreaList{
		\Area{1}
		\Area{3}
		\Area{4}
		\Area{5}
		\Area{6}
	}	
}{
	\EventList{
		\Event{4796}{ }{1}
		\Event{490003}{ }{1}
		\Event{7621}{ }{3}
		\Event{1039126}{ }{3}
		\Event{14349}{ }{4}
		\Event{1052267}{ }{4}
		\Event{12675}{ }{5}
		\Event{1103570}{ }{5}
		\Event{14261}{ }{6}
		\Event{553215}{ }{6}
		
	}
}
\PL{9}{098}{
	\AreaList{
		\Area{3}
		\Area{4}
		\Area{5}
	}	
}{
	\EventList{
		\Event{7324}{o}{3}
		\Event{807892}{ }{3}
		\Event{9676}{o}{3}
		\Event{154411}{ }{3}
		\Event{10288}{o}{4}
		\Event{346984}{ }{4}
		\Event{14349}{ }{4}
		\Event{868129}{ }{4}
		\Event{12675}{ }{5}
		\Event{882168}{ }{5}
		
	}
}
\PL{9}{099}{
	\AreaList{
		\Area{3}
		\Area{4}
		\Area{5}
	}	
}{
	\EventList{
		\Event{7324}{ }{3}
		\Event{1171328}{ }{3}
		\Event{9676}{ }{3}
		\Event{106867}{ }{3}
		\Event{10288}{ }{4}
		\Event{1030259}{ }{4}
		\Event{12675}{ }{5}
		\Event{1544454}{ }{5}
		
	}
}
\PL{9}{100}{
	\AreaList{
		\Area{3}
		\Area{4}
	}	
}{
	\EventList{
		\Event{7324}{ }{3}
		\Event{452779}{ }{3}
		\Event{9676}{ }{3}
		\Event{111314}{ }{3}
		\Event{10288}{ }{4}
		\Event{182910}{ }{4}
		\Event{9653}{o}{4}
		\Event{1921075}{ }{4}
		
	}
}
\PL{9}{101}{
	\AreaList{
		\Area{3}
		\Area{4}
	}	
}{
	\EventList{
		\Event{7324}{ }{3}
		\Event{667662}{ }{3}
		\Event{9676}{ }{3}
		\Event{1018359}{ }{3}
		\Event{10288}{ }{4}
		\Event{834483}{ }{4}
		\Event{9653}{ }{4}
		\Event{170860}{ }{4}
		
	}
}
\PL{9}{102}{
	\AreaList{
		\Area{4}
	}	
}{
	\EventList{
		\Event{9653}{ }{4}
		\Event{166489}{ }{4}
		
	}
}
\PL{9}{103}{
	\AreaList{
		\Area{4}
	}	
}{
	\EventList{
		\Event{9653}{ }{4}
		\Event{1483748}{ }{4}
		
	}
}
\PL{9}{104}{
	\AreaList{
		\Area{2}
		\Area{4}
	}	
}{
	\EventList{
		\Event{10758}{o}{2}
		\Event{116840}{ }{2}
		\Event{14098}{o}{4}
		\Event{1101392}{ }{4}
		
	}
}
\PL{9}{105}{
	\AreaList{
		\Area{2}
		\Area{3}
		\Area{4}
	}	
}{
	\EventList{
		\Event{10758}{ }{2}
		\Event{1002730}{ }{2}
		\Event{1036}{o}{3}
		\Event{120229}{ }{3}
		\Event{14098}{ }{4}
		\Event{1515424}{ }{4}
		
	}
}
\PL{9}{106}{
	\AreaList{
		\Area{2}
		\Area{3}
		\Area{4}
		\Area{5}
		\Area{6}
	}	
}{
	\EventList{
		\Event{10758}{ }{2}
		\Event{1199799}{ }{2}
		\Event{1036}{ }{3}
		\Event{956806}{ }{3}
		\Event{14098}{ }{4}
		\Event{128064}{ }{4}
		\Event{12406}{o}{5}
		\Event{405824}{ }{5}
		\Event{13736}{o}{6}
		\Event{279088}{ }{6}
		
	}
}
\PL{9}{107}{
	\AreaList{
		\Area{2}
		\Area{3}
		\Area{4}
		\Area{5}
		\Area{6}
	}	
}{
	\EventList{
		\Event{10758}{ }{2}
		\Event{491341}{ }{2}
		\Event{1036}{ }{3}
		\Event{958904}{ }{3}
		\Event{2084}{o}{3}
		\Event{943846}{ }{3}
		\Event{14098}{ }{4}
		\Event{753198}{ }{4}
		\Event{12450}{o}{4}
		\Event{163448}{ }{4}
		\Event{12406}{ }{5}
		\Event{398288}{ }{5}
		\Event{13736}{ }{6}
		\Event{249941}{ }{6}
		
	}
}
\PL{9}{108}{
	\AreaList{
		\Area{3}
		\Area{4}
		\Area{5}
		\Area{6}
	}	
}{
	\EventList{
		\Event{1036}{ }{3}
		\Event{547819}{ }{3}
		\Event{2084}{ }{3}
		\Event{41454}{ }{3}
		\Event{12450}{ }{4}
		\Event{191963}{ }{4}
		\Event{12406}{ }{5}
		\Event{626275}{ }{5}
		\Event{13736}{ }{6}
		\Event{674851}{ }{6}
		
	}
}
\PL{9}{109}{
	\AreaList{
		\Area{3}
		\Area{4}
		\Area{5}
		\Area{6}
	}	
}{
	\EventList{
		\Event{2084}{ }{3}
		\Event{224887}{ }{3}
		\Event{12450}{ }{4}
		\Event{1223584}{ }{4}
		\Event{12406}{ }{5}
		\Event{373842}{ }{5}
		\Event{13736}{ }{6}
		\Event{664920}{ }{6}
		
	}
}
\PL{9}{110}{
	\AreaList{
		\Area{1}
		\Area{3}
		\Area{4}
		\Area{5}
	}	
}{
	\EventList{
		\Event{1861}{o}{1}
		\Event{271258}{ }{1}
		\Event{3451}{o}{1}
		\Event{1456560}{ }{1}
		\Event{2084}{ }{3}
		\Event{456626}{ }{3}
		\Event{12450}{ }{4}
		\Event{859715}{ }{4}
		\Event{9814}{o}{5}
		\Event{2118613}{ }{5}
		
	}
}
\PL{9}{111}{
	\AreaList{
		\Area{1}
		\Area{3}
	}	
}{
	\EventList{
		\Event{1861}{ }{1}
		\Event{1768613}{ }{1}
		\Event{3451}{ }{1}
		\Event{990400}{ }{1}
		\Event{9814}{ }{3}
		\Event{1301430}{ }{3}
		
	}
}
\PL{9}{112}{
	\AreaList{
		\Area{1}
		\Area{2}
		\Area{5}
	}	
}{
	\EventList{
		\Event{1861}{ }{1}
		\Event{2268038}{ }{1}
		\Event{3451}{ }{1}
		\Event{2548912}{ }{1}
		\Event{13507}{o}{2}
		\Event{3015226}{ }{2}
		\Event{9814}{ }{5}
		\Event{1311316}{ }{5}
		
	}
}
\PL{9}{113}{
	\AreaList{
		\Area{1}
		\Area{4}
		\Area{5}
	}	
}{
	\EventList{
		\Event{1861}{ }{1}
		\Event{2239182}{ }{1}
		\Event{3451}{ }{1}
		\Event{1412062}{ }{1}
		\Event{9314}{o}{1}
		\Event{2671721}{ }{1}
		\Event{13507}{ }{4}
		\Event{1430942}{ }{4}
		\Event{14353}{o}{4}
		\Event{2671721}{ }{4}
		\Event{9814}{ }{5}
		\Event{828312}{ }{5}
		\Event{11531}{o}{5}
		\Event{2261018}{ }{5}
		
	}
}
\PL{9}{114}{
	\AreaList{
		\Area{1}
		\Area{2}
		\Area{4}
		\Area{5}
	}	
}{
	\EventList{
		\Event{9314}{ }{1}
		\Event{665364}{ }{1}
		\Event{7864}{o}{2}
		\Event{504736}{ }{2}
		\Event{1015}{o}{4}
		\Event{162859}{ }{4}
		\Event{13507}{ }{4}
		\Event{872587}{ }{4}
		\Event{14353}{ }{4}
		\Event{665652}{ }{4}
		\Event{11531}{ }{5}
		\Event{2246502}{ }{5}
		
	}
}
\PL{9}{115}{
	\AreaList{
		\Area{1}
		\Area{2}
		\Area{4}
		\Area{5}
	}	
}{
	\EventList{
		\Event{9314}{ }{1}
		\Event{604205}{ }{1}
		\Event{7864}{ }{2}
		\Event{509552}{ }{2}
		\Event{1015}{ }{4}
		\Event{843686}{ }{4}
		\Event{13507}{ }{4}
		\Event{116643}{ }{4}
		\Event{14353}{ }{4}
		\Event{2050514}{ }{4}
		\Event{11531}{ }{5}
		\Event{950003}{ }{5}
		
	}
}
\PL{9}{116}{
	\AreaList{
		\Area{1}
		\Area{2}
		\Area{3}
		\Area{4}
		\Area{5}
	}	
}{
	\EventList{
		\Event{9314}{ }{1}
		\Event{1471224}{ }{1}
		\Event{15170}{o}{2}
		\Event{2056684}{ }{2}
		\Event{7864}{ }{2}
		\Event{1352363}{ }{2}
		\Event{2709}{o}{3}
		\Event{591614}{ }{3}
		\Event{1015}{ }{4}
		\Event{147195}{ }{4}
		\Event{14353}{ }{4}
		\Event{583376}{ }{4}
		\Event{11531}{ }{5}
		\Event{907693}{ }{5}
		
	}
}
\PL{9}{117}{
	\AreaList{
		\Area{2}
		\Area{3}
		\Area{4}
	}	
}{
	\EventList{
		\Event{15170}{ }{2}
		\Event{549827}{ }{2}
		\Event{7864}{ }{2}
		\Event{1204694}{ }{2}
		\Event{2709}{ }{3}
		\Event{1350417}{ }{3}
		\Event{14870}{o}{3}
		\Event{1073129}{ }{3}
		\Event{1015}{ }{4}
		\Event{1700870}{ }{4}
		\Event{1784}{o}{4}
		\Event{2267683}{ }{4}
		
	}
}
\PL{9}{118}{
	\AreaList{
		\Area{2}
		\Area{3}
		\Area{4}
	}	
}{
	\EventList{
		\Event{15170}{ }{2}
		\Event{503451}{ }{2}
		\Event{14302}{o}{3}
		\Event{808159}{ }{3}
		\Event{2709}{ }{3}
		\Event{578324}{ }{3}
		\Event{14870}{ }{3}
		\Event{1070218}{ }{3}
		\Event{1784}{ }{4}
		\Event{264739}{ }{4}
		
	}
}
\PL{9}{119}{
	\AreaList{
		\Area{2}
		\Area{3}
		\Area{4}
	}	
}{
	\EventList{
		\Event{15170}{ }{2}
		\Event{1468612}{ }{2}
		\Event{14302}{ }{3}
		\Event{1201733}{ }{3}
		\Event{2709}{ }{3}
		\Event{2065600}{ }{3}
		\Event{14870}{ }{3}
		\Event{401360}{ }{3}
		\Event{3679}{o}{3}
		\Event{2813497}{ }{3}
		\Event{1784}{ }{4}
		\Event{1108660}{ }{4}
		\Event{4095}{o}{4}
		\Event{1628345}{ }{4}
		
	}
}
\PL{9}{120}{
	\AreaList{
		\Area{3}
		\Area{4}
	}	
}{
	\EventList{
		\Event{14302}{ }{3}
		\Event{1187514}{ }{3}
		\Event{1658}{o}{3}
		\Event{165626}{ }{3}
		\Event{2368}{o}{3}
		\Event{1927959}{ }{3}
		\Event{14870}{ }{3}
		\Event{1313589}{ }{3}
		\Event{3679}{ }{3}
		\Event{1411386}{ }{3}
		\Event{6556}{o}{4}
		\Event{928624}{ }{4}
		\Event{1784}{ }{4}
		\Event{726785}{ }{4}
		\Event{4095}{ }{4}
		\Event{726785}{ }{4}
		
	}
}
\PL{9}{121}{
	\AreaList{
		\Area{3}
		\Area{4}
	}	
}{
	\EventList{
		\Event{14302}{ }{3}
		\Event{1627677}{ }{3}
		\Event{1658}{ }{3}
		\Event{761221}{ }{3}
		\Event{2368}{ }{3}
		\Event{352170}{ }{3}
		\Event{3679}{ }{3}
		\Event{1026294}{ }{3}
		\Event{6556}{ }{4}
		\Event{662291}{ }{4}
		\Event{4095}{ }{4}
		\Event{62062}{ }{4}
		
	}
}
\PL{9}{122}{
	\AreaList{
		\Area{2}
		\Area{3}
		\Area{4}
		\Area{6}
	}	
}{
	\EventList{
		\Event{14144}{o}{2}
		\Event{1577773}{ }{2}
		\Event{1658}{ }{3}
		\Event{966029}{ }{3}
		\Event{2368}{ }{3}
		\Event{1522407}{ }{3}
		\Event{3679}{ }{3}
		\Event{731737}{ }{3}
		\Event{2377}{o}{4}
		\Event{789487}{ }{4}
		\Event{4272}{o}{4}
		\Event{1871100}{ }{4}
		\Event{6556}{ }{4}
		\Event{252657}{ }{4}
		\Event{4095}{ }{4}
		\Event{924634}{ }{4}
		\Event{6911}{o}{6}
		\Event{421441}{ }{6}
		
	}
}
\PL{9}{123}{
	\AreaList{
		\Area{2}
		\Area{3}
		\Area{4}
		\Area{5}
		\Area{6}
	}	
}{
	\EventList{
		\Event{14144}{ }{2}
		\Event{485325}{ }{2}
		\Event{1658}{ }{3}
		\Event{332942}{ }{3}
		\Event{2368}{ }{3}
		\Event{1452139}{ }{3}
		\Event{2377}{ }{4}
		\Event{317366}{ }{4}
		\Event{4272}{ }{4}
		\Event{526795}{ }{4}
		\Event{6556}{ }{4}
		\Event{1362688}{ }{4}
		\Event{13828}{o}{5}
		\Event{1508211}{ }{5}
		\Event{6911}{ }{6}
		\Event{1033364}{ }{6}
		
	}
}
\PL{9}{124}{
	\AreaList{
		\Area{2}
		\Area{4}
		\Area{5}
		\Area{6}
	}	
}{
	\EventList{
		\Event{14144}{ }{2}
		\Event{916723}{ }{2}
		\Event{14046}{o}{4}
		\Event{697174}{ }{4}
		\Event{2377}{ }{4}
		\Event{301923}{ }{4}
		\Event{4272}{ }{4}
		\Event{1711545}{ }{4}
		\Event{13828}{ }{5}
		\Event{361170}{ }{5}
		\Event{6911}{ }{6}
		\Event{170260}{ }{6}
		
	}
}
\PL{9}{125}{
	\AreaList{
		\Area{2}
		\Area{4}
		\Area{5}
		\Area{6}
	}	
}{
	\EventList{
		\Event{14144}{ }{2}
		\Event{461641}{ }{2}
		\Event{14046}{ }{4}
		\Event{290876}{ }{4}
		\Event{2377}{ }{4}
		\Event{526833}{ }{4}
		\Event{4272}{ }{4}
		\Event{705896}{ }{4}
		\Event{14047}{o}{4}
		\Event{478693}{ }{4}
		\Event{13828}{ }{5}
		\Event{398286}{ }{5}
		\Event{6911}{ }{6}
		\Event{793057}{ }{6}
		
	}
}
\PL{9}{126}{
	\AreaList{
		\Area{3}
		\Area{4}
		\Area{5}
	}	
}{
	\EventList{
		\Event{9297}{o}{3}
		\Event{1454502}{ }{3}
		\Event{14046}{ }{4}
		\Event{1384329}{ }{4}
		\Event{14047}{ }{4}
		\Event{223860}{ }{4}
		\Event{9719}{o}{5}
		\Event{404682}{ }{5}
		\Event{13828}{ }{5}
		\Event{623812}{ }{5}
		
	}
}
\PL{9}{127}{
	\AreaList{
		\Area{3}
		\Area{4}
		\Area{5}
	}	
}{
	\EventList{
		\Event{9297}{ }{3}
		\Event{1873450}{ }{3}
		\Event{14046}{ }{4}
		\Event{38463}{ }{4}
		\Event{14047}{ }{4}
		\Event{897414}{ }{4}
		\Event{9719}{ }{5}
		\Event{1315151}{ }{5}
		
	}
}
\PL{9}{128}{
	\AreaList{
		\Area{2}
		\Area{3}
		\Area{4}
		\Area{5}
	}	
}{
	\EventList{
		\Event{258}{o}{2}
		\Event{1457956}{ }{2}
		\Event{8879}{o}{2}
		\Event{728555}{ }{2}
		\Event{9297}{ }{3}
		\Event{286447}{ }{3}
		\Event{6079}{o}{4}
		\Event{659702}{ }{4}
		\Event{14047}{ }{4}
		\Event{721110}{ }{4}
		\Event{2664}{o}{5}
		\Event{1011065}{ }{5}
		\Event{9719}{ }{5}
		\Event{634161}{ }{5}
		
	}
}
\PL{9}{129}{
	\AreaList{
		\Area{2}
		\Area{3}
		\Area{4}
		\Area{5}
	}	
}{
	\EventList{
		\Event{258}{ }{2}
		\Event{724503}{ }{2}
		\Event{8879}{ }{2}
		\Event{2118329}{ }{2}
		\Event{9297}{ }{3}
		\Event{1805887}{ }{3}
		\Event{6079}{ }{4}
		\Event{439436}{ }{4}
		\Event{13105}{o}{4}
		\Event{840908}{ }{4}
		\Event{2664}{ }{5}
		\Event{2170959}{ }{5}
		\Event{9719}{ }{5}
		\Event{1272572}{ }{5}
		
	}
}
\PL{9}{130}{
	\AreaList{
		\Area{2}
		\Area{4}
		\Area{5}
	}	
}{
	\EventList{
		\Event{258}{ }{2}
		\Event{1042823}{ }{2}
		\Event{8879}{ }{2}
		\Event{552307}{ }{2}
		\Event{508}{o}{4}
		\Event{915589}{ }{4}
		\Event{6079}{ }{4}
		\Event{439885}{ }{4}
		\Event{9657}{o}{4}
		\Event{2248026}{ }{4}
		\Event{13105}{ }{4}
		\Event{175479}{ }{4}
		\Event{2664}{ }{5}
		\Event{795237}{ }{5}
		
	}
}
\PL{9}{131}{
	\AreaList{
		\Area{2}
		\Area{4}
		\Area{5}
	}	
}{
	\EventList{
		\Event{258}{ }{2}
		\Event{897860}{ }{2}
		\Event{8879}{ }{2}
		\Event{930003}{ }{2}
		\Event{508}{ }{4}
		\Event{444676}{ }{4}
		\Event{6079}{ }{4}
		\Event{440097}{ }{4}
		\Event{9657}{ }{4}
		\Event{585689}{ }{4}
		\Event{13105}{ }{4}
		\Event{195477}{ }{4}
		\Event{2664}{ }{5}
		\Event{203440}{ }{5}
		
	}
}
\PL{9}{132}{
	\AreaList{
		\Area{4}
	}	
}{
	\EventList{
		\Event{508}{ }{4}
		\Event{519032}{ }{4}
		\Event{9657}{ }{4}
		\Event{721053}{ }{4}
		\Event{13105}{ }{4}
		\Event{197334}{ }{4}
		
	}
}
\PL{9}{133}{
	\AreaList{
		\Area{4}
	}	
}{
	\EventList{
		\Event{508}{ }{4}
		\Event{929148}{ }{4}
		\Event{9657}{ }{4}
		\Event{2736723}{ }{4}
		
	}
}
%\PL{6}{042}{
	\AreaList{
		\Area{4}
		\Area{6}
	}	
}{
	\EventList{
		\Event{5228}{ }{4}
		\Event{82721}{ }{4}
		\Event{595}{ }{6}
		\Event{375799}{ }{6}
		
	}
}
\PL{6}{043}{
	\AreaList{
		\Area{4}
	}	
}{
	\EventList{
		\Event{5228}{ }{4}
		\Event{82288}{ }{4}
		\Event{8742}{o}{4}
		\Event{531777}{ }{4}
		
	}
}

\PL{6}{045}{
	\AreaList{
		\Area{4}
	}	
}{
	\EventList{
		\Event{8742}{ }{4}
		\Event{288010}{ }{4}
		
	}
}

\PL{6}{078}{
	\AreaList{
		\Area{2}
		\Area{3}
	}	
}{
	\EventList{
		\Event{7535}{ }{2}
		\Event{693350}{ }{2}
		\Event{3615}{ }{3}
		\Event{913}{ }{3}
		\Event{3638}{ }{3}
		\Event{913}{ }{3}
		
	}
}




\end{document}